% R008 — Guide to Cultivating Complex Analysis (Bahasa Indonesia)
% Reader-facing translation unit: source/ca-v1.9/ca.tex lines 1587–1644.
% Source release: v1.9 (commit a4d25d9ba37a486a5a020219eb8bd4e6602fd14c).
% Translation provenance: OpenAI Codex gpt-5.6-sol, Ultra, acting on the user's request.
% Preserve source equations, notation, references, labels, comments, and exercise states.

Dengan demikian, masuk akal untuk menyebut nilai $\nicefrac{1}{z}$ di titik asal
sebagai $\infty$, dan nilai $z$ di $\infty$ sebagai $\infty$.
Bahkan, setiap polinomial yang tidak konstan bernilai $\infty$ di $\infty$.

\begin{exbox}
\begin{exercise}%[\neededexmark]
\label{exercise:polygoesinf}
Misalkan $P(z) = a_d z^d + a_{d-1} z^{d-1} + \cdots + a_1 z + a_0$ adalah
suatu polinomial dengan $a_0,\ldots,a_d \in \C$.  Buktikan bahwa jika $d \geq 1$
dan $a_d \neq 0$ ($P$ tidak konstan),
maka $\lim_{z \to \infty} P(z) = \infty$.
Petunjuk: Jika $a_d = 1$, maka dengan menggunakan
$\abs{\frac{a_{d-1} z^{d-1} + \cdots + a_1 z + a_0}{z^d}}$, diperoleh
$\abs{P(z)} \geq \frac{1}{2} \sabs{z}^d$ untuk $z$ bermodulus cukup besar.
\end{exercise}
\end{exbox}

Dalam kalkulus dan analisis real dasar, Anda mungkin pernah menjumpai
limit tak hingga dalam pengertian bilangan real diperluas.
Meskipun kedua jenis limit tak hingga, baik dalam pengertian
bilangan real diperluas maupun dalam pengertian bola Riemann,
tampak serupa dan pada dasarnya menggunakan notasi yang sama,
keduanya berbeda.
Sebagai contoh,
\begin{equation*}
\lim_{x \to 0} \frac{1}{x} \quad \text{tidak ada,} \qquad \text{tetapi} \qquad
\lim_{z \to 0} \frac{1}{z} = \infty .
\end{equation*}
Di ruas kiri, kita secara tersirat menggunakan pengertian bilangan real diperluas
($x$ real, bukan?) dan di
ruas kanan, kita secara tersirat menggunakan pengertian bola Riemann ($z$ tampaknya kompleks).
Asumsi tersirat semacam itu tentu saja dapat menimbulkan kebingungan.
Dalam buku ini, limit akan dipahami dalam pengertian bola Riemann
kecuali dinyatakan lain atau sudah jelas.
Jika mungkin timbul kebingungan, kita akan menulis $\infty$ pada bilangan real diperluas
sebagai $+\infty$.

\medskip

Aritmetika (parsial) yang dapat kita
definisikan secara wajar dengan $\infty$ pada bola Riemann
cukup berbeda dari
$\infty$ pada bilangan real diperluas.
Untuk $c \neq \infty$, kita dapat mendefinisikan $c+\infty = \infty$,
tetapi baik $\infty+\infty$ maupun $\infty-\infty$ tidak bermakna.
Sebagai contoh, jika $f(z) = z$ dan $g(z) = c-z$,
maka $\lim_{z \to \infty} f(z) = \infty$,
$\lim_{z \to \infty} g(z) = \infty$, tetapi $\lim_{z \to \infty}
\bigl(f(z)+g(z)\bigr) = c$, sehingga $\infty+\infty$ tidak bermakna.
Berbeda dari bilangan real diperluas,
wajar (dan terkadang berguna) untuk mendefinisikan $\nicefrac{c}{0} = \infty$ untuk
$c \neq 0$, $\nicefrac{c}{\infty} = 0$ untuk $c \neq \infty$,
dan $c \cdot \infty = \infty$ untuk $c \neq 0$.
%We could also define $\infty \cdot \infty = \infty$, $\frac{0}{\infty} = 0$,
%and
%$\frac{\infty}{0} = \infty$.
Sebaiknya ungkapan $0 \cdot \infty$, $\nicefrac{0}{0}$, dan
$\nicefrac{\infty}{\infty}$ dibiarkan tidak terdefinisi.
